\chapter{INTRODUCCIÓN}
\label{ch:intro}
 
El Sistema de Cuentas Nacionales (SCN) constituye el marco estadístico internacional de referencia para medir la actividad económica de un país. Dentro de este marco, existe un conjunto de organizaciones cuya lógica no responde a la maximización del lucro ni al control estatal, y cuya contribución no se refleja de forma agregada en los principales indicadores macroeconómicos: las instituciones sin fines de lucro (ISFL) y, en un sentido más amplio, las organizaciones de la economía social. Reconociendo esta limitación, Naciones Unidas publicó el \textit{Handbook on Non-profit Institutions in the System of National Accounts} y, posteriormente, el \textit{Satellite Account on Non-profit and Related Institutions and Volunteer Work} (2018), que constituye el capítulo específico del SCN dedicado a medir el aporte de este universo de organizaciones al producto interno bruto \cite{vereinte_nationen_satellite_2018}. Este marco reconoce a las ISFL, las cooperativas, las mutuales y otras entidades afines como parte de un sector con dinámica económica propia, cuya visibilización estadística resulta condición necesaria para el diseño de políticas públicas informadas.\\
 
Dentro de este universo, las cooperativas constituyen una de las formas organizacionales más difundidas de la economía social a nivel mundial. La literatura y los manuales metodológicos internacionales reconocen distintas ramas de acuerdo con su actividad y el vínculo con sus socios: cooperativas de trabajo asociado, cooperativas agrícolas y campesinas, cooperativas de consumo, cooperativas de vivienda, cooperativas de servicios y cooperativas financieras --entre las que se incluyen las mutuales de seguros y las cooperativas de ahorro y crédito-- \cite{barea_tejeiro_manual_2007}. Cada una de estas ramas presenta particularidades en cuanto a su forma de operación, su regulación y su tratamiento contable dentro del SCN, pero todas comparten los principios cooperativos de propiedad democrática, participación de los socios y reinversión de excedentes en la organización o la comunidad \shortcite{perilleux_are_2016}. En Chile, este universo cooperativo abarca cooperativas agrícolas, de vivienda, de trabajo, de consumo y de servicios, entre otras, cuya diversidad ha sido documentada por la División de Asociatividad y Cooperativas (DAES) y por la literatura académica reciente \shortcite{radrigan_rubio_organizaciones_2024}.\\
 
En este panorama, las cooperativas de ahorro y crédito (CAC) constituyen un actor especialmente relevante del sistema financiero chileno. A diferencia de la banca, estas organizaciones operan bajo un modelo de propiedad democrática: los usuarios son simultáneamente socios, las decisiones se adoptan bajo el principio de un socio, un voto, y los excedentes se reinvierten en la organización o se distribuyen entre los miembros según su participación \cite{vereinte_nationen_satellite_2018}. Esta estructura les permite operar en territorios y segmentos de la población donde la banca convencional tiene escasa presencia, movilizando el ahorro local y facilitando el acceso al crédito para pequeños emprendedores y trabajadores \shortcite{perilleux_are_2016}. \\
 
En Chile, las CAC se encuentran reguladas principalmente por el Decreto con Fuerza de Ley N°5 de 2003 y la Ley General de Cooperativas \cite{dfl5_cooperativas_2003}. Su supervisión es dual: la División de Asociatividad y Cooperativas (DAES) del Ministerio de Economía ejerce la supervisión de gobierno corporativo sobre la totalidad del sector, mientras que la Comisión para el Mercado Financiero (CMF) regula a aquellas CAC con patrimonio superior a las 400.000 UF. Con la promulgación de la Ley N°21.641 de Resiliencia Financiera en diciembre de 2023, la CMF asumió además la supervisión
prudencial integral de las cooperativas de mayor tamaño
\cite{ley21641_resiliencia_2023}. Pese a su relevancia económica y social, no existe en Chile una medición sistemática del aporte de estas organizaciones al producto interno bruto, ni una comparación con otros actores del sistema financiero, como la banca comercial o el mercado de valores, lo que limita tanto la formulación de políticas públicas como la evaluación de su impacto en el desarrollo económico del país.\\
 
Ante este vacío, la experiencia internacional ofrece una metodología consolidada: las cuentas satélite de la economía social. Una cuenta satélite es un instrumento estadístico que, bajo el marco metodológico del Sistema de Cuentas Nacionales (SCN), permite medir el aporte económico de sectores que no quedan adecuadamente captados en las estadísticas nacionales estándar \cite{noauthor_cuenta_2006}. En el ámbito de la economía social, Naciones Unidas publicó en 2018 el \textit{Satellite Account on Non-profit and Related Institutions and
Volunteer Work}, que establece el marco metodológico internacional para construir este tipo de cuentas, incorporando explícitamente a las cooperativas elegibles dentro de su alcance \cite{vereinte_nationen_satellite_2018}. A nivel europeo, el manual elaborado por CIRIEC para la Comisión Europea ofrece una guía específica para cooperativas y mutuas, detallando cómo traducir su contabilidad interna a las variables del SCN para estimar el Valor Agregado Bruto (VAB) del sector \cite{barea_tejeiro_manual_2007}.\\
 
Diversos países han avanzado en la implementación de esta metodología. Francia fue pionera en desarrollar cuentas satélite para la economía social desde la década de 1980, seguida por Bélgica, que publicó sus primeras cuentas en 2004. España, por su parte, inició en 2024 la construcción formal de su Cuenta Satélite de la Economía Social bajo la conducción del Instituto Nacional de Estadística, aplicando directamente el manual de Naciones Unidas de 2018 \cite{noauthor_cuenta_2024}. Estos ejercicios han permitido a dichos países cuantificar la relevancia del sector cooperativo y de la economía social en el producto interno bruto, visibilizar su contribución al empleo y orientar políticas públicas específicas para el sector. Estas experiencias demuestran que la metodología es técnicamente viable y que sus resultados son comparables entre países \cite{archambault_lester_2022}.
 
Chile cuenta con los elementos necesarios para aplicar esta metodología. Por un lado, existe experiencia previa en la construcción de cuentas satélite sectoriales: la Cuenta Satélite de Tecnologías de Información y Comunicación, desarrollada por el Ministerio de Economía en conjunto con el INE, demostró que es posible integrar registros administrativos dispersos para estimar el aporte de un sector al PIB nacional \cite{noauthor_cuenta_2006}. Por otro lado, los registros administrativos necesarios existen: la DAES mantiene el catastro de cooperativas activas, la CMF publica información financiera de las CAC
bajo su supervisión, y el SII dispone de registros administrativos de actividad económica para las personas jurídicas activas. El desafío no es la ausencia de datos, sino su fragmentación y la falta de un sistema que los articule bajo un marco metodológico común \shortcite{radrigan_rubio_organizaciones_2024}.
 
En este contexto, la presente memoria propone estimar el aporte de las
cooperativas de ahorro y crédito al PIB chileno mediante la metodología
de cuentas satélite, adoptando el SCN 2025 como marco de consistencia
macroeconómica y el Manual ONU-TSE \shortcite{vereinte_nationen_satellite_2018}
como referencia metodológica principal, e integrando los registros
administrativos disponibles en Chile a través de un prototipo de panel interactivo de visualización (plataforma Huella Social) que consolida información proveniente de diversas fuentes institucionales y permite generar indicadores comparables y replicables para el sector de la economía social \shortcite{munoz_exploring_2021}. La hipótesis
central que guía este trabajo es que Chile dispone de los registros
administrativos suficientes para construir una estimación del Valor Agregado
Bruto (VAB) de las CAC, sentando así las bases para una futura Cuenta
Satélite de la Economía Social en el país.\footnote{Específicamente, los registros administrativos utilizados en esta memoria incluyen el archivo maestro de personas jurídicas del SII (\texttt{PUB\_NOMBRES\_PJ.txt}) y la serie de tiempo (\texttt{sii\_company\_timeseries.parquet}), que reportan tramo de ventas y número de trabajadores para las organizaciones registradas.}
 
\newpage

\section{Marco institucional del proyecto}
 
Esta memoria se desarrolla en el marco del proyecto \textbf{Huella Social},
una iniciativa de investigación aplicada orientada a construir un sistema
integrado de información para organizaciones de la economía social en Chile.
El proyecto surge de la colaboración entre la Universidad de los Andes y el
Instituto Nacional de Asociatividad y Cooperativas (INAC), organismo
dependiente del Ministerio de Economía, Fomento y Turismo.
 
El trabajo aquí presentado constituye uno de los principales productos de
formación de capital humano del proyecto, y sus resultados metodológicos
y conceptuales sirvieron como insumo directo para la postulación de
\textit{Huella Social} al concurso \textbf{FONDEF IDeA I+D 2026} de ANID.
El detalle de la trayectoria previa del proyecto, sus resultados
tecnológicos comprometidos y sus socios estratégicos se presenta en el
Anexo~\ref{anx:fondef}.
 


\section{Objetivos}
 
\subsection{Objetivo general}
 
Estimar el aporte de las cooperativas de ahorro y crédito al
producto interno bruto de Chile mediante la metodología de
cuentas satélite, integrando registros administrativos
disponibles en una plataforma conceptual de datos que permita generar indicadores comparables y replicables para el sector de la economía social.
 
\subsection{Objetivos específicos}
 
\begin{enumerate}
 
\item Analizar el estado del arte y la literatura relevante sobre economía social, transparencia institucional y sistemas de información para organizaciones sociales.
 
\item Identificar y caracterizar las principales fuentes de información disponibles sobre cooperativas en Chile.
 
\item Estimar el aporte al PIB nacional de las cooperativas de ahorro y crédito en Chile mediante la construcción de una cuenta satélite sectorial para el período 2014--2024.
 
\item Desarrollar un panel interactivo de visualización que integre los resultados de la cuenta satélite y permita consultar indicadores de desempeño financiero e institucional de las cooperativas de ahorro y crédito.

\end{enumerate}
 Los objetivos específicos siguen una secuencia lógica: el primero sitúa el
problema en la literatura especializada y construye el aparato conceptual
necesario (Capítulos~\ref{ch:revision} y~\ref{ch:marco}); el segundo
identifica y sistematiza las fuentes de información administrativas
disponibles en Chile (Sección~\ref{sec:fuentes-informacion} del
Capítulo~\ref{ch:metodologia}), insumo indispensable para el tercero, que
aplica dicha información a la construcción de la cuenta satélite y la
estimación del aporte al PIB (Sección~\ref{sec:metodos-estimacion} del
Capítulo~\ref{ch:metodologia} y Capítulo~\ref{ch:resultados}); finalmente,
el cuarto objetivo traduce estos resultados en una herramienta de
visualización, cuyo diseño se presenta en la
Sección~\ref{sec:panel} del Capítulo~\ref{ch:metodologia} y cuyas
visualizaciones se reportan en la Sección~\ref{sec:dashboard} del
Capítulo~\ref{ch:resultados}.
\section{Alcances}
 
La presente memoria se enfoca en la estimación del aporte de las cooperativas de ahorro y crédito (CAC) al PIB nacional mediante la construcción de una cuenta satélite sectorial para el período 2014--2024. El trabajo considera el análisis de las fuentes de información disponibles, la revisión de literatura académica sobre economía social y cuentas satélite, y la aplicación de la metodología del Sistema de Cuentas Nacionales 2025.
 
El alcance del estudio se limita a las CAC activas con información financiera disponible en los registros de la DAES y la CMF. 

Esta delimitación no constituye una restricción del enfoque, sino la estrategia de pilotaje definida en la postulación FONDEF del proyecto Huella Social (Anexo A.1): el formulario establece a las cooperativas de ahorro y crédito como caso piloto de la plataforma de cuentas satélite, con escalamiento posterior al conjunto de la economía social y solidaria. La presente memoria ejecuta ese piloto y avanza sobre lo declarado en el formulario, transitando desde el diseño conceptual comprometido en la postulación hacia una estimación verificable del aporte al PIB para un tipo específico de organización. La extensión de la metodología a otros tipos de cooperativas y organizaciones queda fuera del alcance de este trabajo, y corresponde a la agenda del proyecto FONDEF descrita en el Anexo ~\ref{anx:resultados tecnologicos}.
 
\section{Estructura del documento}
 
 
El presente trabajo se organiza en seis capítulos, complementados con un conjunto de anexos que sustentan el desarrollo metodológico y proveen información de contexto.
 
El Capítulo~\ref{ch:intro} corresponde a la presente introducción, y entrega el marco institucional del proyecto, los objetivos generales y específicos, los alcances y la estructura del documento.
 
El Capítulo~\ref{ch:revision} presenta la revisión de literatura, abordando los antecedentes relevantes sobre economía social, cooperativas de ahorro y crédito, transparencia institucional, sistemas de información para organizaciones sociales, y experiencias internacionales en la construcción de cuentas satélite.
 
El Capítulo~\ref{ch:marco} desarrolla el marco teórico, profundizando en los fundamentos del Sistema de Cuentas Nacionales, la metodología de cuentas satélite y su aplicación al sector cooperativo.
 
El Capítulo~\ref{ch:metodologia} describe la metodología del estudio. Se presenta el marco conceptual del SCN aplicado a las cooperativas de ahorro y crédito, el inventario de fuentes de información disponibles en Chile, los requerimientos establecidos por los manuales internacionales y el diagnóstico de brechas resultante. A continuación, se detallan los métodos de estimación utilizados —integración de fuentes, operacionalización de variables del SCN y modelo de cálculo del Valor Agregado Bruto— y, finalmente, el diseño conceptual de la plataforma Huella Social como aplicación operacional del marco metodológico.
 
El Capítulo~\ref{ch:resultados} presenta los resultados del estudio, incluyendo la estimación del aporte de las cooperativas de ahorro y crédito al PIB, la propuesta de diseño de la plataforma y el análisis de su potencial aplicación.
 
Finalmente, el Capítulo~\ref{ch:conclusion} expone las conclusiones del trabajo, las principales recomendaciones y la agenda de mejoras pendientes para una versión más completa de la cuenta satélite.



